\documentclass[11pt]{article}
\usepackage{amsmath,amssymb}
\title{A common-gauge mixed-grid current estimate with original relative boundaries}\author{}\date{}
\begin{document}\maketitle
\paragraph{Original grids and precisely allowed connectors.}
Use the original scalar grids of the mixed boundary magnetic loop bound calculation: one Neumann coordinate and
three Dirichlet coordinates, or four Dirichlet coordinates, with the
literal nearest-neighbor killed/reflected kernels $p_t^{(\mu)}$.
Their staggered origins and side lengths are retained. A connector is a
fixed translation between overlapping subsets of two grids, of bounded
length, or a composition of a bounded number of such maps.
Its matrix is a partial permutation: row and column sums are at most one.
The reverse connector is its adjoint. A finite sum of the original local
vertex choices changes constants only by their fixed multiplicity.
The current along a connector is the line integral of the SAME constant
Landau field along a fixed bounded axis path. No transition probability
is replaced by a new bath law.

All component scalar phases are first put into the same ambient Landau
gauge $a_1=b x_0$, $b=2/R^2$, by the exact component gauge gradients of
the mixed boundary magnetic loop bound calculation. The corresponding endpoint phases are retained in the
connectors. Paths can have different component grids; their actual
Neumann coordinates need not coincide.

\paragraph{1. A common-gauge open-current moment.}
Fix an ambient starting point $x$ and use the common centered gauge
$a_1=b(z_0-x_0)$. On one grid, let
\[
 J_t=\int_0^t b(Z_{s-}-x_0)\,dY_s,\qquad Z=X^0,\quad Y=X^1,
\]
on surviving paths. If coordinate one is Dirichlet, the stopped-current
argument of the mixed boundary magnetic loop bound calculation gives $E[J_t^2;\tau>t]\le Cb^2t^2$.

If coordinate one is Neumann, coordinate zero MUST be Dirichlet.
Discrete integration by parts (there are no simultaneous coordinate
jumps) gives, on each surviving path,
\[
 J_t=-\int_0^t b(Y_{s-}-x_1)\,dZ_s
                  +b(Z_t-x_0)(Y_t-x_1).                         \tag{1}
\]
The first term, stopped INCLUDING the possible killing jump, is a
martingale with second moment at most $2b^2t^2$.
For the endpoint term extend the coordinate processes past killing.
They are independent; the reflected coordinate is the same Lipschitz
fold used in the mixed boundary magnetic loop bound calculation, and each displacement second moment is at
most $2t$. Dropping the survival event gives
\[
 E[(Z_t-x_0)^2(Y_t-x_1)^2;\tau>t]\le4t^2.
\]
Thus in the SAME gauge for EVERY component grid,
\[
 \boxed{E_x[J_t^2;\tau>t]\le12b^2t^2.}                            \tag{2}
\]
The endpoint product in (1) is retained, not silently declared a
martingale. Equation (2) is an unnormalized moment, not a bridge estimate.

\paragraph{2. Piecewise paths and reversed mixed loops.}
Consider a path consisting of a fixed bounded number of original
semigroup segments, interspersed with the allowed connectors.
Its total positive duration is $t$.
The common-gauge current on a segment equals its current centered at
that segment's start, plus
$b(z_{\rm start}-x_0)\Delta y$.
Conditional displacement moments on each next segment are bounded by
$2t_j$ uniformly in its starting state; connector translations add
bounded deterministic displacements. Repeated conditioning, (2), and
Cauchy--Schwarz for a fixed number of summands therefore give
\[
 E[J_{\rm piecewise}^2;\text{survival}]
                      \le Cb^2(1+t)^2.                           \tag{3}
\]
For example the change-of-center term has second moment at most
$Cb^2(t_{\rm previous}+1)t_j$. A bounded connector contributes at most
$Cb^2(t_{\rm previous}+1)$. These estimates use the actual substochastic
path measures, not independence across conditioned bridges.

Each piecewise kernel is still substochastic in BOTH rows and columns.
Among a bounded number of time segments one has duration at least a
fixed fraction of $t$. Bounding that factor by its original heat-kernel
supremum and summing all other factors proves
\[
 \sup K_{\rm piecewise}\le C k(t),\qquad
 k(t)=(1+t)^{-2}e^{-ct/R^2}.                                      \tag{4}
\]
Constants absorb the bounded segment count and connector offsets.

Now concatenate a path on one grid with the reverse of a path on a
different grid, including their endpoint connectors. Cut the resulting
loop after half its total positive duration and reverse the second
half. Both halves start at the SAME ambient anchor; the reversed
connectors are adjoints. Their currents subtract in the same centered
gauge. Bounding the other half by (4), using (3), and summing the
$O(R^4)$ initial anchors gives the replacement for a same-grid
semigroup loop identity:
\[
 \boxed{\mathcal A_{\rm mixed\,loop}(t)
       \le C R^4 b^2(1+t)^2 k(t).}                               \tag{5}
\]
Reversibility is used separately on each original grid. No fictitious
single-grid semigroup is inserted. The $1+t$ accommodates bounded
connector areas at zero positive time.

\paragraph{3. Three different component grids.}
Take two oppositely charged and one neutral path between two common
ambient vertex anchors, each on any original component or ghost grid.
Their bounded endpoint connectors make the currents comparable.
Assign the neutral path its geometric current only as an auxiliary
variable. This changes no path weight.
If the neutral path is shortest, use
$J_1-J_2=(J_1-J_3)+(J_3-J_2)$.
Otherwise use the two charged paths directly.
Bound the unused path kernel by its original supremum and apply (5)
to each pair. For sorted times $m\le M\le T$ this proves
\[
 \mathcal A_{\theta,\rm mixed}
 \le CR^4b^2(1+T)^2 k(M)k(T).                                   \tag{6}
\]
This is the mixed-grid version of the theta network joint magnetic current calculation; it also handles
arbitrarily unbalanced times. Consequently
\[
 \int\frac{\mathcal A_{\theta,\rm mixed}}{1+T}\,dt_1dt_2dt_3
               \le CR^4b^2[\log(2+R)]^2.                        \tag{7}
\]
For small times the harmless added constant is integrable.
For large times the ordered integration is the same as the theta network joint magnetic current calculation.
The factor $1/(1+T)$ in (7) remains EXPLICITLY UNPAID.

\paragraph{Actual-model scope and the remaining join.}
These estimates use the original component/ghost scalar kernels and
bounded local connector choices. They cover the geometric current of
their charge-conserving path networks even when the reflected
coordinates differ. A field-dependent arbitrary connector, a long
tree path, or a nonlocal projection is outside the stated lemma.
The actual Wilson curl and ghost finite-difference vertices must still
supply the derivative weight, with their differentiated endpoint,
spin and gauge-projection partners retained. The full coefficient also
requires the finite-boundary mass and longitudinal cancellations;
(6) does not replace them.

A finite-state regression checks (2) on every starting point of the
$R=3$ grids with Neumann coordinate zero, one or two and on the all-
Dirichlet ghost grid. It confirms the common-gauge moment including
the case where the current coordinate reflects. This regression is
not the proof of (5)--(7).
No bound on the complete nonlinear coefficient, its original remainder,
or S94 is claimed. The mixed-grid current comparison is the closed
join; the current-weighted derivative estimate remains open.
\end{document}
